\pagestyle{empty}

\begin{center}
	\begin{tikzpicture}[remember picture,overlay]
		\fill[black] (current page.south west) rectangle (current page.north east);

		% Autor
		\node[white, font=\Huge, anchor=north, text width =\linewidth, align = center] (author) at ([yshift=-44pt] current page.north) {\textcolor{customgold}{\scshape\large\authorname}};

		% Título
		\node[white, font=\Large, anchor=north, text width =0.9\linewidth, align = center, below=4pt of author.south] (title) {\scshape\Huge\booktitle};

		% Padrão de círculos central
		\begin{scope}[shift={(current page.center)}, scale=0.75]
			\foreach \i in {0, 45, ..., 315} {
					\fill[customgold] (\i:5) circle (0.5);
				}

			\foreach \i in {5, 4, 3, 2, 1} {
					\fill[customgold] (0,0) circle (\i);
					\fill[white] (0,0) circle (\i-0.25);
					\fill[black] (0,0) circle (\i-0.5);
				}
		\end{scope}



		\node[white, font=\Huge, anchor=south] (publisher) at ([yshift=32pt] current page.south) {\textcolor{gray}{\scshape\large LEM}};


		% Logotipo da editora
		\node[anchor=north, above=4pt of publisher.north] {\includegraphics[width=0.15\textwidth]{figs/LEM.eps}};

	\end{tikzpicture}
\end{center}

\cleardoublepage

\begin{titlepage}
	\centering

	\thispagestyle{empty}

	\vspace*{0.5cm}

	{\scshape\Large \authorname \par}

	\vfill

	\vspace{1.5cm}

	{\scshape\huge \booktitle \par}

	{\vspace{6pt}
		{\scshape\large \subtitle\par}

		\vspace{\stretch{1.25}}}

	\vspace{1.5cm}

	\vfill

	{\scshape\large \editionyear \par}

\end{titlepage}


